% 可数集（综述）
% license CCBYSA3
% type Wiki

本文根据 CC-BY-SA 协议转载翻译自维基百科\href{https://en.wikipedia.org/wiki/Countable_set}{相关文章}

在数学中，如果一个集合是有限的，或者它可以与自然数集合建立一一对应关系，那么它就是可数的。[a] 等价地说，如果存在一个从该集合到自然数的单射函数，那么该集合就是可数的；这意味着集合中的每个元素都可以与一个唯一的自然数对应，或者说集合的元素可以一个接一个地数出来，尽管由于元素个数无限，计数过程可能永远不会结束。

从更技术性的角度讲，假设接受**可数选择公理**，若一个集合的基数（即集合中的元素个数）不大于自然数集合的基数，则它是可数的。一个可数但不是有限的集合称为**可数无穷**。

这一概念归功于格奥尔格·康托尔（Georg Cantor），他证明了**不可数集**的存在，也就是说存在一些集合不是可数的；例如实数集合就是不可数的。

关于术语的说明

尽管这里所定义的“可数”（countable）与“可数无穷”（countably infinite）用法相当普遍，但并非所有文献都采用这种术语。\[1] 一种替代用法是：用 **countable** 来表示这里所说的“可数无穷”，而用 **at most countable** 来表示这里所说的“可数”。\[2]\[3]

术语 **enumerable**\[4] 和 **denumerable**\[5]\[6] 也可能被使用，例如分别指代“可数”和“可数无穷”。\[7] 但其定义并不统一，并且需要注意与“递归可枚举”的区别。\[8]

